% FST-NS_LogDistanceIntegrability_v1_de.tex
\documentclass[12pt,a4paper]{article}

\usepackage[utf8]{inputenc}
\usepackage[T1]{fontenc}
\usepackage{lmodern}
\usepackage{amsmath,amssymb,amsthm}
\usepackage{enumitem}
\usepackage{xcolor}
\usepackage{geometry}
\geometry{a4paper, margin=2.5cm}
\usepackage{hyperref}

\hypersetup{
    colorlinks=true,
    linkcolor=blue!70!black,
    citecolor=green!50!black,
    urlcolor=blue!70!black
}

\theoremstyle{definition}
\newtheorem{definition}{Definition}[section]
\theoremstyle{plain}
\newtheorem{theorem}[definition]{Satz}
\newtheorem{lemma}[definition]{Lemma}
\newtheorem{proposition}[definition]{Proposition}
\newtheorem{corollary}[definition]{Korollar}
\theoremstyle{remark}
\newtheorem{remark}[definition]{Bemerkung}
\newtheorem{example}[definition]{Beispiel}

\renewcommand{\proofname}{Beweis}

\newcommand{\Var}{\mathrm{Var}}
\newcommand{\dist}{\mathrm{dist}}
\newcommand{\AC}{\mathrm{AC}}
\newcommand{\BV}{\mathrm{BV}}
\newcommand{\id}{\mathrm{id}}

\title{\textbf{Log-Distanz-Integrierbarkeit fuer dissipative Stroemungen:\\
BV-Selektionen auf fraktalen Attraktoren}}
\author{Lukas Geiger\\
\small Unabhaengiger Forscher, Bernau im Schwarzwald, Deutschland\\
\small ORCID: \href{https://orcid.org/0009-0005-7296-1534}{0009-0005-7296-1534}%
\thanks{KI-Werkzeuge (Claude von Anthropic) wurden für mathematische Formalisierung, Literaturverifikation und redaktionelle Verfeinerung eingesetzt. Alle mathematischen Inhalte wurden vom Autor entwickelt, verifiziert und freigegeben.}}
\date{Maerz 2026}

\begin{document}
\maketitle

\begin{center}
\small
\textbf{MSC 2020:} 37L30, 47H09, 46B20, 35Q30, 26A45\\
\textbf{Schluesselwoerter:} Naechstpunkt-Projektion, log-Lipschitz-Regularitaet,
beschraenkte Variation, Attraktorgeometrie, dissipative PDEs, Navier--Stokes
\end{center}

\begin{abstract}
Wir fuehren einen Regularitaetsrahmen fuer Naechstpunkt-Projektionen
auf kompakte Attraktoren dissipativer Semiflows in Hilbertraeumen ein.
Das klassische Lipschitz-Projektionstheorem (Federer, 1959) erfordert
positive Reichweite -- eine Bedingung, die mit fraktaler
Attraktorgeometrie unvereinbar ist.  Wir ersetzen sie durch eine
\emph{trajektorienweise log-Lipschitz}-Bedingung (TLL), die den
Projektionsmodul auf der geometrisch natuerlichen Skala des
Attraktorabstands statt der Partitionsfeinheit kontrolliert.
In Kombination mit einer \emph{Log-Distanz-Integrierbarkeits}-Bedingung
(LDI) ergibt dies Selektionen beschraenkter Variation entlang absolut
stetiger Trajektorien -- genau die Regularitaet, die fuer
Gronwall-artige Argumente in bedingten Regularitaetskriterien
benoetigt wird.  Der Rahmen gilt fuer jeden dissipativen Semiflow
mit kompaktem Attraktor endlicher fraktaler Dimension, einschliesslich
der 3D-Navier-Stokes-Gleichungen, Reaktions-Diffusions-Systeme und
gedaempfter Wellengleichungen.  Zwei verifizierbare hinreichende
Bedingungen fuer LDI werden identifiziert: beschraenkte Variation
der Log-Distanz-Funktion und exponentielle Sublevel-Zeitmass-Abnahme.
\end{abstract}


% ==================================================================
\section{Einleitung}
% ==================================================================

Sei $H$ ein separabler Hilbertraum und $\mathcal{A} \subset H$ eine
kompakte Menge -- der globale Attraktor eines dissipativen Semiflows
$\{S(t)\}_{t \ge 0}$.  Eine grundlegende Frage in der qualitativen
Theorie dissipativer partieller Differentialgleichungen ist die
Regularitaet der Naechstpunkt-Projektion
\[
  P(x) \;\in\; \arg\min_{a \in \mathcal{A}} \|x - a\|_H
\]
als Funktion der Zeit entlang einer Trajektorie $u(t) = S(t)\,u_0$.
Ist die Abbildung $t \mapsto P(u(t))$ absolut stetig (AC), so
schliessen sich Gronwall-artige Energieabschaetzungen automatisch;
ist sie lediglich messbar, bleiben kritische Fehlerterme
unkontrolliert.

Die klassische Theorie (Federer~\cite{Federer1959}) garantiert
Lipschitz-Stetigkeit der Naechstpunkt-Projektion auf der tubularen
Umgebung von Mengen mit \emph{positiver Reichweite}.  Attraktoren
dissipativer PDEs haben jedoch generisch nicht-ganzzahlige fraktale
Dimension und Reichweite null
(vgl.~\cite{Temam1997}, Kapitel~V).  Dies schliesst die direkte
Anwendung von Federers Theorem aus.

Ziel dieser Arbeit ist die Einfuehrung einer schwaecheren
Regularitaetsklasse -- \emph{trajektorienweise log-Lipschitz}-Projektionen
-- die mit fraktaler Geometrie vertraeglich ist und dennoch fuer die
Kontrolle der projizierten Trajektorie durch beschraenkte Variation (BV)
genuegt.  Motiviert durch log-Lipschitz-Regularitaetsergebnisse fuer
dissipative Stroemungen
(vgl.~\cite{BarbuCannone2016}) beobachten wir, dass die naive globale
log-Lipschitz-Bedingung
\begin{equation}\label{eq:global-LL-intro}
  \|P(x) - P(y)\|
  \;\le\;
  C\,\|x-y\|\,\bigl(1 + \log\tfrac{R}{\|x-y\|}\bigr)
\end{equation}
\emph{nicht} BV von $P \circ u$ fuer AC-Kurven~$u$ impliziert: Die
Partitionssummen $\sum_i \Delta_i(1+\log(R/\Delta_i))$ divergieren
logarithmisch bei Verfeinerung des Gitters
(Abschnitt~\ref{sec:counterexample}).  Die korrekte Formulierung
ersetzt das Inkrementskalen-Gewicht $\log(R/\|x-y\|)$ durch das
Attraktorabstands-Gewicht $\log(R/d(t))$, wobei
$d(t) = \dist(u(t),\mathcal{A})$.  Dies liefert eine
trajektorienadaptierte Bedingung, deren BV-Summen konvergieren,
wann immer eine Log-Distanz-Integrierbarkeitsbedingung (LDI) erfuellt ist.


% ==================================================================
\section{Grundlagen und Notation}
% ==================================================================

Durchgehend bezeichnet $H$ einen separablen reellen Hilbertraum mit
Norm $\|\cdot\|$ und innerem Produkt $\langle\cdot,\cdot\rangle$.

\begin{definition}[Dissipativer Semiflow und Attraktor]
\label{def:semiflow}
Ein \emph{dissipativer Semiflow} auf $H$ ist eine Familie
$\{S(t)\}_{t \ge 0}$ stetiger Abbildungen $S(t): H \to H$ mit
$S(0) = \id$, $S(t+s) = S(t) \circ S(s)$, die einen
\emph{kompakten globalen Attraktor}
$\mathcal{A} \subset H$ besitzt: eine kompakte, invariante
($S(t)\mathcal{A} = \mathcal{A}$ fuer alle $t \ge 0$) Menge, die
jede beschraenkte Teilmenge von~$H$ anzieht.
\end{definition}

\begin{definition}[Tubulare Umgebung und Abstand]
\label{def:tube}
Fuer $R > 0$ ist die \emph{$R$-tubulare Umgebung} von $\mathcal{A}$
\[
  \mathcal{A}^R
  \;:=\;
  \{x \in H : \dist(x, \mathcal{A}) < R\},
\]
wobei $\dist(x, \mathcal{A}) := \inf_{a \in \mathcal{A}} \|x - a\|$.
Fuer eine Kurve $u: [0,T] \to \mathcal{A}^R$ schreiben wir
$d(t) := \dist(u(t), \mathcal{A})$.
\end{definition}

\begin{definition}[Naechstpunkt-Selektion]
\label{def:selection}
Eine \emph{Naechstpunkt-Selektion} ist eine messbare Abbildung
$P: \mathcal{A}^R \to \mathcal{A}$ mit
$\|x - P(x)\| = \dist(x, \mathcal{A})$ fuer alle
$x \in \mathcal{A}^R$ und $P|_{\mathcal{A}} = \id$.
Eine solche Selektion existiert nach dem Satz von
Kuratowski--Ryll-Nardzewski, angewandt auf die kompaktwertige
Multifunktion
$x \mapsto \arg\min_{a \in \mathcal{A}} \|x-a\|$.
\end{definition}

\begin{proposition}[Minimale metrische Ableitungsselektion]
\label{prop:minimal-selection}
Unter den Naechstpunkt-Selektionen
$P: \mathcal{A}^R \to \mathcal{A}$ existiert eine
\emph{kanonische Selektion} $u^*(t)$ entlang jeder absolut stetigen
Trajektorie $u \in \AC([0,T];\,H)$, definiert durch
\[
  u^*(t)
  \;:=\;
  \arg\min_{a \in \mathcal{A},\;\|u(t)-a\|=d(t)}
  |a|'(t),
  \qquad
  |a|'(t)
  \;:=\;
  \lim_{h \to 0}
  \frac{\dist(a,\,P(u(t+h)))}{|h|},
\]
die die metrische Ableitung
$|u^*|'(t) = \lim_{h\to 0} \|u^*(t+h) - u^*(t)\|/|h|$
unter allen Naechstpunkt-Selektionen minimiert. Diese Selektion ist
f.ue.\ eindeutig: Fuer f.ue.\ $t \in [0,T]$ ist der Minimierer der
metrischen Ableitung ueber der (kompakten) Naechstpunkt-Menge
$\{a \in \mathcal{A} : \|u(t)-a\| = d(t)\}$ eindeutig.
\end{proposition}

\begin{proof}
Die Naechstpunkt-Menge $\Pi(t) := \{a \in \mathcal{A} :
\|u(t)-a\| = d(t)\}$ ist kompakt und nichtleer fuer jedes~$t$.
Das Funktional der metrischen Ableitung $a \mapsto
\limsup_{h\to 0}\|P_{a}(u(t+h)) - a\|/|h|$ ist unterhalbstetig auf
der kompakten Menge~$\Pi(t)$, nimmt also sein Minimum an. Fuer
f.ue.~$t$, an dem $|u'|(t)$ existiert, hat die Naechstpunkt-Menge
ein eindeutiges Element, das die metrische Ableitung minimiert: Falls
zwei verschiedene Minimierer $a_1, a_2 \in \Pi(t)$ auf einer Menge
positiven Masses existierten, wuerde die Konvexitaet der Norm
$\|(a_1+a_2)/2 - u(t)\| < d(t)$ fuer den Mittelpunkt liefern
(strikte Konvexitaet der Hilbertnorm), im Widerspruch zu
$a_1,a_2 \in \mathcal{A}$. Also gilt Eindeutigkeit f.ue.
\end{proof}

\begin{remark}
Die kanonische Selektion $u^*$ ist der natuerliche Kandidat fuer TLL
(Definition~\ref{def:TLL}): Unter allen messbaren
Naechstpunkt-Selektionen hat sie die kleinste metrische Ableitung
und damit die besten Chancen, die log-Lipschitz-Schranke zu erfuellen.
Im Abschnitt Einschraenkungen haben wir festgestellt, dass verschiedene
Selektionen unterschiedliche Regularitaet haben koennen; die minimale
metrische Ableitungsselektion ist die ,,beste''.
\end{remark}


% ==================================================================
\section{Warum globale Log-Lipschitz-Stetigkeit versagt}
\label{sec:counterexample}
% ==================================================================

Wir zeigen zunaechst, dass die globale log-Lipschitz-Bedingung
\eqref{eq:global-LL-intro}, selbst kombiniert mit
Log-Distanz-Integrierbarkeit, nicht direkt BV von $P \circ u$ liefert.

\begin{proposition}[Versagen der globalen LL-Schranke zur Variationskontrolle]
\label{prop:divergence}
Sei $P: \mathcal{A}^R \to \mathcal{A}$ mit
\eqref{eq:global-LL-intro} und Konstante $C > 0$.  Sei
$u: [0,1] \to \mathcal{A}^R$ Lipschitz mit
$\|u'(t)\| = 1$ f.ue.  Fuer die gleichmaessige Partition
$t_i = i/n$ ($i = 0, \ldots, n$) erfuellt die aus
\eqref{eq:global-LL-intro} gewonnene obere Schranke
\[
  \sum_{i=0}^{n-1} \|P(u(t_{i+1})) - P(u(t_i))\|
  \;\le\;
  C\,(1 + \log(nR)),
\]
die fuer $n \to \infty$ divergiert.  Insbesondere liefert die
globale LL-Bedingung keine endliche \emph{a priori}-Schranke
fuer $\Var(P \circ u)$.
\end{proposition}

\begin{proof}
Jedes Inkrement erfuellt $\Delta_i := \|u(t_{i+1}) - u(t_i)\| = 1/n$.
Nach \eqref{eq:global-LL-intro}:
\[
  \sum_{i=0}^{n-1} \|P(u(t_{i+1})) - P(u(t_i))\|
  \;\le\;
  C \sum_{i=0}^{n-1} \frac{1}{n}\,
  \bigl(1 + \log(nR)\bigr)
  \;=\;
  C\,(1 + \log(nR)).
\]
Fuer $n \to \infty$ waechst dies wie $C\log n \to \infty$.
Die aus \eqref{eq:global-LL-intro} abgeleitete Schranke divergiert
also mit der Gitterverfeinerung und liefert keine endliche obere
Schranke fuer $\Var(P \circ u)$.  (Ob die Variation selbst unendlich
ist, haengt von der spezifischen Geometrie von~$\mathcal{A}$ ab;
entscheidend ist, dass die globale LL-Bedingung als \emph{Werkzeug}
zur Herleitung von BV unzureichend ist.)
\end{proof}

\begin{remark}
Die Divergenz ist \emph{logarithmisch}, nicht polynomiell -- viel
langsamer als fuer H\"older-Abbildungen (wo $\sum \Delta_i^\alpha$
fuer $\alpha < 1$ konvergiert).  Dies macht globales LL ,,beinahe''
hinreichend fuer BV, aber die logarithmische Akkumulation ueber
unendlich viele Partitionspunkte verhindert es.  Die trajektorienweise
log-Lipschitz-Bedingung vermeidet dies, indem sie das Log-Gewicht am
Attraktorabstand $d(t)$ verankert, der nicht von der Partition
abhaengt.
\end{remark}


% ==================================================================
\section{Die trajektorienweise log-Lipschitz-Bedingung}
\label{sec:TLL}
% ==================================================================

\begin{definition}[Trajektorienweise log-Lipschitz-Projektion (TLL)]
\label{def:TLL}
Sei $u \in \AC([0,T];\, H)$ mit $u(t) \in \mathcal{A}^R$ fuer
alle~$t$, und sei $P: \mathcal{A}^R \to \mathcal{A}$ eine
Naechstpunkt-Selektion.  Wir sagen, dass $P$ die
\emph{trajektorienweise log-Lipschitz-Bedingung} (TLL) entlang~$u$
erfuellt, wenn Konstanten $C_{\mathrm{TLL}} > 0$ und $h_0 > 0$
existieren, sodass fuer alle $t \in [0,T]$ mit $d(t) > 0$ und alle
$h \in \mathbb{R}$ mit $0 < |h| < h_0$ und
$t+h \in [0,T]$
(d.h.\ sowohl Vorwaerts- als auch Rueckwaerts-Inkremente) gilt:
\begin{equation}\label{eq:TLL}
  \|P(u(t+h)) - P(u(t))\|
  \;\le\;
  C_{\mathrm{TLL}}\,\|u(t+h) - u(t)\|\,
  \Omega(t),
  \tag{TLL}
\end{equation}
wobei $\Omega(t) := 1 + \log(R/d(t))$ das
\emph{Log-Distanz-Gewicht} ist.

Auf der Menge $\{t : d(t) = 0\}$ (wo $u(t) \in \mathcal{A}$)
gilt $P(u(t)) = u(t)$ durch Idempotenz von~$P$.  Fuer den
Uebergang $d(t) = 0$, $d(t+h) > 0$ gilt
$\|P(u(t+h)) - u(t)\| = \|P(u(t+h)) - P(u(t))\|$.
Da $u(t) \in \mathcal{A}$ ein Kandidat fuer den naechsten Punkt
von $u(t+h)$ ist, liefert die Minimalitaet von $P$:
$\|u(t+h) - P(u(t+h))\| \le \|u(t+h) - u(t)\|$.
Zusammen mit der Dreiecksungleichung:
\[
  \|P(u(t+h)) - u(t)\|
  \;\le\;
  \|P(u(t+h)) - u(t+h)\| + \|u(t+h) - u(t)\|
  \;\le\;
  2\,\|u(t+h) - u(t)\|.
\]
Dies ergibt $|v'|(t) \le 2\,|u'|(t)$ auf $\{d = 0\}$.
\end{definition}

\begin{remark}[Konvention $C_{\mathrm{TLL}} \ge 2$]
\label{rem:CTLL-convention}
Der Faktor~$2$ auf $\{d = 0\}$ wird durch die TLL-Schranke
absorbiert, sofern $C_{\mathrm{TLL}} \ge 2$ und $\Omega(t) \ge 1$.
Durchgehend verwenden wir die Konvention
$C_{\mathrm{TLL}} := \max(C_{\mathrm{TLL}},\,2)$;
dies aendert die BV-Schranke \eqref{eq:BV-bound} um hoechstens
einen Faktor~$2$ und beeinflusst die Integrierbarkeit
von~\eqref{eq:LDI} nicht.
\end{remark}

\begin{remark}[Geometrische Interpretation von TLL]
\label{rem:TLL-geometry}
Die Bedingung TLL besagt, dass die Naechstpunkt-Projektion,
ausgewertet entlang der Trajektorie~$u$, eine lokale
Lipschitz-Konstante hat, die \emph{hoechstens logarithmisch}
divergiert, wenn sich die Trajektorie dem Attraktor naehert:
\[
  \mathrm{Lip}_{\mathrm{loc}}(P \circ u;\, t)
  \;\lesssim\;
  1 + \log\frac{R}{d(t)}.
\]
Im Abstand $d(t) \sim R$ vom Attraktor ist die Projektion
im Wesentlichen Lipschitz (der Log-Term ist $O(1)$).  Fuer
$d(t) \to 0$ waechst das Log-Gewicht und spiegelt die zunehmende
geometrische Komplexitaet der Attraktorgrenze auf feinen Skalen
wider.  Die logarithmische Rate ist die \emph{kritische Schwelle}:
Schnelleres Wachstum (z.B. polynomiell) wuerde LDI generisch
zum Scheitern bringen, waehrend langsameres Wachstum (beschraenkt)
Lipschitz-Projektion implizieren wuerde -- zu stark fuer fraktale
Geometrie.
\end{remark}

\begin{remark}[TLL als Aubin-Eigenschaft der Projektionskorrespondenz]
\label{rem:TLL-Aubin}
Die TLL-Bedingung laesst eine natuerliche Umformulierung in der
Sprache der Variationsanalyse zu.  Definiere die
\emph{Projektionskorrespondenz}
\[
  \Gamma: t \;\mapsto\;
  \bigl\{a \in \mathcal{A} : \|a - u(t)\| = d(t)\bigr\},
\]
die jeder Zeit die Menge der naechsten Attraktorelemente zuordnet.
Die TLL-Bedingung $|d'(t)| \le C/|\log d(t)|$
(d.h., die metrische Ableitung der Abstandsfunktion wird durch den
inversen Logarithmus kontrolliert) ist aequivalent zur
\emph{Aubin-Eigenschaft} (metrische Regularitaet) der Korrespondenz
$\Gamma$ mit Rate $1/|\log r|$.  Konkret ist $\Gamma$
metrisch regulaer mit Rate $\phi(r) = 1/|\log r|$, wenn
\[
  \dist\bigl(a,\, \Gamma(t')\bigr)
  \;\le\;
  L\,\phi\bigl(\|u(t') - u(t)\|\bigr)\,\|u(t') - u(t)\|
\]
fuer $a \in \Gamma(t)$ und $t'$ nahe~$t$ gilt.
Diese Umformulierung verbindet die Projektionsregularitaetstheorie
mit dem Aubin-Kriterium und dem Rahmenwerk metrischer Regularitaet
von Rockafellar--Wets~\cite{RockafellarWets1998} und
Dontchev--Rockafellar~\cite{DontchevRockafellar2009}.
Der zentrale Vorteil ist, dass Aubin-Eigenschafts-Kriterien fuer
mengenwertige Abbildungen gut entwickelt sind und verifizierbare
hinreichende Bedingungen liefern (z.B.\ ueber die
Mordukhovich-Koderivierte), die auf die Attraktorprojektion
anwendbar sein koennten.
\end{remark}

\begin{remark}[TLL vs.\ globales LL]
\label{rem:TLL-vs-LL}
Die globale log-Lipschitz-Bedingung \eqref{eq:global-LL-intro}
impliziert \emph{nicht} TLL.  Fuer $h \to 0$ bei festem~$t$:
\[
  \log\frac{R}{\|u(t+h)-u(t)\|}
  \;\approx\;
  \log\frac{R}{h\,|u'(t)|}
  \;=\;
  \log\frac{R}{|u'(t)|} + \log\frac{1}{h}
  \;\to\; \infty,
\]
waehrend $\log(R/d(t))$ fest bleibt.  Die globale Bedingung
erzeugt ein \emph{Inkrementskalen}-Log-Gewicht, das mit der
Partitionsfeinheit divergiert; TLL erzeugt ein
\emph{Abstandsskalen}-Log-Gewicht, das nur von der aktuellen
Geometrie abhaengt.  Dies ist der strukturelle Grund, warum TLL
BV liefert, waehrend globales LL dies nicht tut.
\end{remark}


% ==================================================================
\section{Log-Distanz-Integrierbarkeit}
\label{sec:LDI}
% ==================================================================

\begin{definition}[Log-Distanz-Integrierbarkeit (LDI)]
\label{def:LDI}
Im Rahmen von Definition~\ref{def:TLL} sagen wir, dass
\emph{Log-Distanz-Integrierbarkeit} gilt, wenn
\begin{equation}\label{eq:LDI}
  \int_{\{d>0\}} \|u'(t)\|\,\Omega(t)\,\mathrm{d}t
  \;=\;
  \int_{\{d>0\}} \|u'(t)\|\,
  \bigl(1 + \log\tfrac{R}{d(t)}\bigr)\,\mathrm{d}t
  \;<\; \infty.
  \tag{LDI}
\end{equation}
Auf der Menge $\{t : d(t) = 0\}$ gilt $v(t) = u(t)$ und der
BV-Beitrag wird durch $\int_{\{d=0\}} \|u'(t)\|\,\mathrm{d}t
< \infty$ kontrolliert (da $u$ AC ist), sodass nur die Region
$\{d > 0\}$ fuer \eqref{eq:LDI} relevant ist.  Um die
BV-Schranke \eqref{eq:BV-bound} einheitlich auf ganz $[0,T]$
zu formulieren, setzen wir $\Omega(t) := 1$ auf $\{d=0\}$
(konsistent mit der Faktor-$2$-Schranke aus
Bemerkung~\ref{rem:CTLL-convention}).
\end{definition}

\begin{remark}[Minimalitaet von LDI]
\label{rem:LDI-minimal}
Die Bedingung LDI ist praezise zwischen den bekannten
Regularitaetsklassen positioniert:
\begin{itemize}
  \item \textbf{Lipschitz-Projektion} ($\Omega \equiv 1$):
    LDI reduziert sich auf $\int \|u'\| < \infty$, was fuer
    AC-Kurven automatisch gilt.  Dies ist das Federer-/positive-Reichweite-Regime.
  \item \textbf{H\"older-Projektion} ($\Omega(t) \sim d(t)^{-\beta}$,
    $\beta > 0$): LDI erfordert $\int \|u'\|\,d(t)^{-\beta} < \infty$,
    was generisch in der Naehe des Attraktors scheitert.  Deshalb
    liefern Man\'e-Projektoren (H\"older-Inverse) kein BV.
  \item \textbf{Log-Lipschitz-Projektion} (TLL):
    LDI erfordert $\int \|u'\|\,\log(1/d) < \infty$, den
    Grenzfall -- logarithmische Singularitaeten sind unter milden
    dynamischen Bedingungen integrierbar.
\end{itemize}
Das log-Lipschitz/LDI-Paar besetzt somit das \emph{kritische}
Regularitaetsfenster fuer BV-Selektionen auf fraktalen Attraktoren.
\end{remark}


% ==================================================================
\section{Hauptresultat: BV-Kettenregel}
\label{sec:main}
% ==================================================================

\begin{theorem}[BV-Kettenregel fuer log-Lipschitz-Projektionen]
\label{thm:BV-chain}
Sei $\{S(t)\}_{t \ge 0}$ ein dissipativer Semiflow auf einem
separablen Hilbertraum~$H$ mit kompaktem globalem Attraktor
$\mathcal{A} \subset H$.  Sei $u \in \AC([0,T];\, H)$ mit
$u(t) \in \mathcal{A}^R$ fuer alle~$t$, und sei
$P: \mathcal{A}^R \to \mathcal{A}$ eine Naechstpunkt-Selektion,
die TLL \eqref{eq:TLL} entlang~$u$ erfuellt.  Setze $v(t) := P(u(t))$.

Wenn LDI \eqref{eq:LDI} gilt, dann ist
$v \in \BV([0,T];\, H)$ mit der expliziten Schranke
\begin{equation}\label{eq:BV-bound}
  \Var_H(v;\,[0,T])
  \;\le\;
  C_{\mathrm{TLL}}
  \int_0^T \|u'(t)\|\,\Omega(t)\,\mathrm{d}t,
\end{equation}
wobei $\Omega(t) := 1 + \log(R/d(t))$ auf $\{d > 0\}$ und
$\Omega(t) := 1$ auf $\{d = 0\}$
(Bemerkung~\ref{rem:CTLL-convention}).
\end{theorem}

\begin{proof}
Der Beweis gliedert sich in drei Schritte.

\medskip
\noindent
\textbf{Schritt 1 (Schranke der metrischen Ableitung).}\;
Fuer jede stetige Kurve $v: [0,T] \to H$ ist die \emph{obere
metrische Ableitung} bei $t$ definiert als
\[
  |v'|^+(t)
  \;:=\;
  \limsup_{h \to 0} \frac{\|v(t+h) - v(t)\|}{|h|}
  \;\in\; [0,+\infty].
\]
Diese Groesse existiert ueberall (als Wert in $[0,+\infty]$)
und ist Borel-messbar.  Fuer absolut stetige Kurven ist der
$\limsup$ f.ue.\ ein genuiner Limes
(vgl.~\cite[Theorem~1.1.2]{AmbrosioGigliSavare2008}).
Da $u$ AC ist, existiert $|u'|(t) := \lim_{h\to 0}\|u(t+h)-u(t)\|/|h|$
fuer f.ue.~$t$.

\smallskip\noindent
\emph{Fall $d(t)>0$:}\;
Division von \eqref{eq:TLL} durch $|h|$ und Bildung des $\limsup$:
\begin{equation}\label{eq:metric-deriv-bound}
  |v'|^+(t)
  \;\le\;
  C_{\mathrm{TLL}}\,|u'|(t)\,\Omega(t)
  \qquad \text{fuer f.ue.\ } t \text{ mit } d(t) > 0.
\end{equation}
(Die rechte Seite ist ein genuiner Limes; die linke Seite kann
ein $\limsup$ sein, der durch dieselbe Schranke dominiert wird.)

\smallskip\noindent
\emph{Fall $d(t)=0$ (Inneres von $\mathcal{Z}$):}\;
Auf $\mathcal{Z} := \{t \in [0,T] : d(t) = 0\}$ gilt
$u(t) \in \mathcal{A}$ und $v(t) = P(u(t)) = u(t)$.
Da $u(t) \in \mathcal{A}$ ein Kandidat-Minimierer fuer
$\dist(u(t+h), \mathcal{A})$ ist, liefert die
Dreiecksungleichung
$\|P(u(t+h)) - u(t)\| \le 2\,\|u(t+h) - u(t)\|$
(vgl.\ die Schranke in Definition~\ref{def:TLL}),
also $|v'|^+(t) \le 2\,|u'|(t)$.

\smallskip\noindent
\emph{Randpunkte von $\mathcal{Z}$:}\;
An $t_0 \in \partial\mathcal{Z}$ mit $d(t_0)=0$ koennen Folgen
$t_n \to t_0$ mit $d(t_n) > 0$ existieren.  Fuer die Stetigkeit von
$v$ bei $t_0$: $\|v(t_n) - v(t_0)\| = \|P(u(t_n)) - u(t_0)\|
\le d(t_n) + \|u(t_n) - u(t_0)\| \to 0$, da $u$
stetig ist und $d = \dist(\cdot,\mathcal{A}) \circ u$ Stetigkeit
erbt (die Abstandsfunktion ist $1$-Lipschitz).  Die Schranke der metrischen Ableitung
$|v'|^+(t_0) \le 2\,|u'|(t_0)$ bleibt durch dasselbe
Dreiecksungleichungsargument gueltig.

\smallskip\noindent
Da $\Omega(t) \ge 1$ und $C_{\mathrm{TLL}} \ge 2$
(Bemerkung~\ref{rem:CTLL-convention}), gilt die
Schranke \eqref{eq:metric-deriv-bound} auf ganz $[0,T]$
(mit $|v'|^+$ statt $|v'|$, wo der Limes moeglicherweise
nicht existiert; dies genuegt fuer das BV-Kriterium in Schritt~3).

\medskip
\noindent
\textbf{Schritt 2 (Stetigkeit von~$v$).}\;
Die Naechstpunkt-Projektion auf eine kompakte (aber nichtkonvexe)
Menge ist im Allgemeinen nur oberhalbstetig als mengenwertige
Abbildung und muss nicht ueberall einwertig sein.  Die
TLL-Bedingung \eqref{eq:TLL} impliziert jedoch direkt die Stetigkeit
von $v = P \circ u$ entlang der Trajektorie: Fuer $d(t) > 0$ gilt
$\|v(t+h) - v(t)\| \le C_{\mathrm{TLL}}\,\|u(t+h)-u(t)\|\,
\Omega(t) \to 0$ fuer $h \to 0$; fuer $d(t) = 0$ gilt
$\|v(t+h)-v(t)\| \le 2\,\|u(t+h)-u(t)\| \to 0$
(einschliesslich Randpunkte von $\mathcal{Z}$, vgl.\ Schritt~1).
Also ist $v$ stetig auf $[0,T]$.

\medskip
\noindent
\textbf{Schritt 3 (BV-Schranke).}\;
Wir zeigen $\|v(b) - v(a)\| \le \int_a^b g(t)\,\mathrm{d}t$
fuer alle $0 \le a < b \le T$, wobei
$g(t) := C_{\mathrm{TLL}}\,|u'|(t)\,\Omega(t) \in L^1$
(nach LDI).  Da $v$ stetig ist (Schritt~2), folgt die
BV-Schranke $\Var(v) \le \int_0^T g$ durch Summation ueber
jede Partition.

Setze $f(t) := \|v(t) - v(a)\|$.  Dann ist $f: [a,b] \to \mathbb{R}$
stetig und reellwertig,
$f(a) = 0$, und die obere Dini-Ableitung erfuellt
$D^+ f(t) \le |v'|^+(t)$ (da
$f(t+h) - f(t) \le \|v(t+h)-v(t)\|$ nach der umgekehrten
Dreiecksungleichung, und Division durch $h > 0$ mit
$\limsup$ ergibt $D^+ f \le |v'|^+$).
Nach Schritt~1 gilt $|v'|^+(t) \le g(t)$ f.ue.  Ein klassisches Resultat ueber Dini-Ableitungen
(vgl.\ \cite[Kapitel~VII, \S\S\,9--10]{Saks1937}; siehe
auch \cite[Kapitel~3]{AmbrosioFuscoPallara2000}) besagt:
Fuer eine stetige Funktion $f:[a,b] \to \mathbb{R}$ mit
$D^+ f(t) \le g(t)$ f.ue.\ und $g \in L^1$ gilt die
Ungleichung $f(b) - f(a) \le \int_a^b g(t)\,\mathrm{d}t$.
Daher
\[
  \|v(b) - v(a)\| \;=\; f(b)
  \;\le\; \int_a^b g(t)\,\mathrm{d}t
  \;\le\; \int_0^T g(t)\,\mathrm{d}t.
\]
Summation ueber jede Partition $0 = t_0 < \cdots < t_n = T$:
\[
  \sum_i \|v(t_{i+1}) - v(t_i)\|
  \;\le\;
  \sum_i \int_{t_i}^{t_{i+1}} g(t)\,\mathrm{d}t
  \;=\;
  \int_0^T g(t)\,\mathrm{d}t
  \;=\;
  C_{\mathrm{TLL}} \int_0^T |u'|(t)\,\Omega(t)\,\mathrm{d}t,
\]
was \eqref{eq:BV-bound} ergibt.

\medskip
\noindent
Die Schranke \eqref{eq:BV-bound} folgt aus den Schritten~1--3.
\end{proof}

\begin{remark}[Direkte Verifikation ueber Partitionssummen]
\label{rem:partition-sum}
Fuer eine feste Partition mit Feinheit $< h_0$ laesst sich die
Endlichkeit der Partitionssumme direkt aus TLL ablesen, ohne die
Dini-Ableitungs-Maschinerie: Auf jedem Teilintervall liefert TLL
$\|v(t_{i+1}) - v(t_i)\|
\le C_{\mathrm{TLL}}\,\|u(t_{i+1}) - u(t_i)\|\,\Omega(t_i)$
(oder die Faktor-$2$-Schranke auf $\{d = 0\}$, absorbiert durch
$C_{\mathrm{TLL}} \ge 2$).  Die resultierende Partitionssumme
$\sum_i C_{\mathrm{TLL}}\,\Omega(t_i)\int_{t_i}^{t_{i+1}}
\|u'\|\,\mathrm{d}t$ ist fuer jede feste Partition endlich
(da $\Omega$ auf jeder kompakten Menge disjunkt von $\{d=0\}$
beschraenkt ist und der $\{d=0\}$-Beitrag durch $2\int\|u'\|$
kontrolliert wird).  Der Uebergang zum Supremum ueber \emph{alle}
Partitionen -- also zur Definition von
$\Var(v)$ -- erfordert jedoch das Dini-Ableitungs-Argument
von Schritt~3, das das punktweise $\Omega(t_i)$-Gewicht
zum integrierten $\Omega(t)$-Gewicht verschaerft.
\end{remark}

\begin{remark}[Banachraum-Allgemeinheit]
\label{rem:banach}
Der Beweis von Satz~\ref{thm:BV-chain} verwendet nur die Norm,
Kompaktheit von~$\mathcal{A}$ und messbare Selektion -- kein
inneres Produkt und keine Hilbertraum-Struktur.  Die
BV-Kettenregel gilt daher wortwortlich in jedem separablen
Banachraum.  Wir beschraenken uns nur wegen der PDE-Anwendungen
in Abschnitt~\ref{sec:applications} auf Hilbertraeume.
\end{remark}


% ==================================================================
\section{Hinreichende Bedingungen fuer LDI}
\label{sec:sufficient}
% ==================================================================

Die Bedingung LDI \eqref{eq:LDI} ist eine einzelne
Integrierbarkeitsanforderung.  Wir geben zwei unabhaengige
hinreichende Bedingungen an, die jeweils aus dynamischen
Informationen ueber den Semiflow verifizierbar sind.

% ------------------------------------------------------------------
\subsection{Variante A: Log-Distanz hat beschraenkte Variation}
% ------------------------------------------------------------------

\begin{proposition}[BV-Log-Distanz impliziert LDI]
\label{prop:variant-A}
Angenommen $d(t) > 0$ fuer f.ue.\ $t \in [0,T]$ und
$\log(R/d(t)) \in \BV([0,T])$.  Dann gilt LDI.
\end{proposition}

\begin{proof}
Schreibe $\ell(t) := 1 + \log(R/d(t))$.  Da $u' \in L^1(0,T;\,H)$
(weil $u$ AC ist) und $\ell \in \BV([0,T]) \subset
L^\infty([0,T])$, ist das Produkt
$\|u'(t)\|\,\ell(t)$ lokal integrierbar:
\[
  \int_0^T \|u'(t)\|\,\ell(t)\,\mathrm{d}t
  \;\le\;
  \|\ell\|_{L^\infty(0,T)}\,\int_0^T \|u'(t)\|\,\mathrm{d}t
  \;<\; \infty.  \qedhere
\]
\end{proof}

\begin{remark}[Wann ist die Log-Distanz BV?]
\label{rem:log-BV}
Die Bedingung $\log(R/d) \in \BV$ erfordert, dass die Trajektorie
nicht unendlich oft zwischen verschiedenen Abstandsskalen zum
Attraktor oszilliert.  Formal schliesst sie unendliche
,,Zickzack-Frequenz'' in $d(t)$ aus.  Fuer monoton anziehende
Trajektorien ($d(t)$ schliesslich fallend) ist dies automatisch
erfuellt.  Fuer Trajektorien mit transienten Oszillationen folgt
BV von $\log d$, wenn die Gesamtzahl der Skalenkreuzungen
$\{t : d(t) = \varepsilon\}$ fuer jedes
$\varepsilon > 0$ endlich ist -- eine generische Eigenschaft
dissipativer Semiflows mit isolierten Gleichgewichten.
\end{remark}

% ------------------------------------------------------------------
\subsection{Variante B: Sublevel-Zeitmass-Kontrolle}
% ------------------------------------------------------------------

\begin{definition}[Sublevel-Zeit-Kontrolle (STC)]
\label{def:STC}
Die Trajektorie~$u$ erfuellt \emph{Sublevel-Zeit-Kontrolle} mit
Exponent $\alpha > 0$, wenn ein $C_* > 0$ existiert mit
\begin{equation}\label{eq:STC}
  \bigl|\{t \in [0,T] : d(t) \le \varepsilon\}\bigr|
  \;\le\;
  C_*\,\varepsilon^\alpha
  \qquad \text{fuer alle } \varepsilon \in (0, R].
  \tag{STC}
\end{equation}
\end{definition}

\begin{remark}[STC impliziert $|\{d=0\}| = 0$]
\label{rem:STC-d0}
Der Grenzuebergang $\varepsilon \to 0$ in \eqref{eq:STC} zeigt
$|\{d = 0\}| = 0$.  STC ist daher nur auf Trajektorien anwendbar,
die keine positive Zeit auf dem Attraktor verbringen.  Dies ist der
generische Fall fuer nicht-stationaere Loesungen, die sich
$\mathcal{A}$ asymptotisch naehern.  Fuer Trajektorien, die den
Attraktor in endlicher Zeit erreichen (z.B.\ Konvergenz zu einem
Gleichgewicht mit $u(t) \in \mathcal{A}$ fuer $t \ge t_1$), ist
Variante~A (BV-Log-Distanz) die geeignete hinreichende Bedingung.
\end{remark}

\begin{proposition}[STC impliziert LDI]
\label{prop:variant-B}
Falls STC \eqref{eq:STC} mit einem $\alpha > 0$ gilt und
$u' \in L^2(0,T;\,H)$, dann ist
$\log(R/d) \in L^p(0,T)$ fuer jedes $p < \infty$, und LDI gilt.
\end{proposition}

\begin{proof}
\textbf{Schritt 1 (Tonelli-Darstellung).}\;
Da $u(t) \in \mathcal{A}^R$, gilt $d(t) \le R$ fuer alle~$t$.
Auf der Menge $\{t : d(t) = 0\}$ (wo $u(t) \in \mathcal{A}$) ist
der Integrand $\|u'(t)\|\,\log(R/d(t)) = +\infty$; wir setzen den
Beitrag per Konvention auf null (auf dieser Menge gilt $v(t) = u(t)$
und die BV-Kettenregel-Schranke ist trivial erfuellt).
Auf $\{t : d(t) > 0\}$ ergibt die Schichtzerlegungsidentitaet
$\log(R/d(t)) = \int_{d(t)}^{R} s^{-1}\,\mathrm{d}s$ mit Tonelli:
\begin{equation}\label{eq:tonelli}
  \int_{\{d>0\}} \|u'(t)\|\,\log\frac{R}{d(t)}\,\mathrm{d}t
  \;=\;
  \int_0^R \frac{1}{s}
  \Bigl(\int_{\{0 < d(t) \le s\}} \|u'(t)\|\,\mathrm{d}t\Bigr)
  \,\mathrm{d}s.
\end{equation}

\medskip
\noindent
\textbf{Schritt 2 (H\"older auf dem inneren Integral).}\;
Nach der H\"olderschen Ungleichung:
\[
  \int_{\{0<d(t)\le s\}} \|u'(t)\|\,\mathrm{d}t
  \;\le\;
  \Bigl(\int_0^T \|u'(t)\|^2\,\mathrm{d}t\Bigr)^{1/2}
  \cdot
  \bigl|\{0 < d \le s\}\bigr|^{1/2}
  \;\le\;
  \|u'\|_{L^2}\,C_*^{1/2}\,s^{\alpha/2},
\]
wobei der letzte Schritt $|\{0<d\le s\}| \le |\{d\le s\}|
\le C_*\,s^\alpha$ aus STC verwendet.

\medskip
\noindent
\textbf{Schritt 3 (Integrierbarkeit nahe $s=0$).}\;
Einsetzen in \eqref{eq:tonelli}:
\[
  \int_0^R \frac{1}{s}\,\|u'\|_{L^2}\,C_*^{1/2}\,s^{\alpha/2}
  \,\mathrm{d}s
  \;=\;
  \|u'\|_{L^2}\,C_*^{1/2}
  \int_0^R s^{\alpha/2 - 1}\,\mathrm{d}s
  \;=\;
  \|u'\|_{L^2}\,C_*^{1/2}\,
  \frac{2}{\alpha}\,R^{\alpha/2}
  \;<\; \infty.
\]
Das Integral konvergiert, weil $\alpha > 0$, sodass
$s^{\alpha/2-1}$ nahe $s=0$ integrierbar ist.

\medskip
\noindent
\textbf{Schritt 4 ($L^p$-Schranken).}\;
Dasselbe Tonelli-Argument mit H\"older-Exponenten $p$ und $q$
($1/p + 1/q = 1$) ergibt
$\|\log(R/d)\|_{L^p} \le C(p,\alpha,R,C_*)$ fuer jedes
$p < \infty$, sodass LDI aus jeder endlichen $L^q$-Schranke
fuer~$\|u'\|$ folgt.
\end{proof}

\begin{remark}[Quellen von STC in dissipativen Stroemungen]
\label{rem:STC-sources}
Die Sublevel-Zeit-Kontrolle \eqref{eq:STC} ist der genuine neue
analytische Input.  Drei Mechanismen koennen sie erzeugen:

\begin{enumerate}[label=(\roman*)]
  \item \textbf{Exponentielles Tracking:}\;
    Falls $d(t) \le C_0\,e^{-\lambda t}$ fuer $t \ge t_0$, dann gilt
    $\{d \le \varepsilon\} \subset
    \{t \ge (\log(C_0/\varepsilon))/\lambda\}$,
    dessen Mass $O(\log(1/\varepsilon))$ ist -- viel kleiner als
    $C_*\,\varepsilon^\alpha$ fuer jedes $\alpha > 0$.
    Exponentielles Tracking impliziert also STC mit \emph{jedem}
    Exponenten.
  \item \textbf{Lyapunov-Dissipation:}\;
    Falls ein Lyapunov-Funktional $E$ die Bedingung
    $-\dot{E} \ge \Psi(d)$ mit nicht zu flachem $\Psi$ nahe $d=0$
    erfuellt (z.B.\ $\Psi(d) \ge c\,d^\beta$ fuer ein $\beta$),
    dann werden die Sublevel-Zeiten durch das gesamte Energiebudget
    $E(0) - \inf E$ kontrolliert.
  \item \textbf{Keine schnellen Oszillationen:}\;
    Energetische Einschraenkungen (endliches Dissipationsintegral)
    verhindern beliebig schnelle Hin-und-Her-Bewegungen nahe dem
    Attraktor und begrenzen die Anzahl der Exkursionen in
    $\{d \le \varepsilon\}$.
\end{enumerate}
Jeder Mechanismus liefert STC mit verschiedenen Exponenten~$\alpha$.
Die STC-Bedingung ist physikalisch plausibel fuer alle
Standardklassen dissipativer PDEs.
\end{remark}

\begin{lemma}[Exponentielles Tracking impliziert Log-Integrierbarkeit]
\label{lem:stc-tracking}
Falls $d(t) \le C\,e^{-\lambda t}$ fuer ein $\lambda > 0$ und
$C > 0$, dann gilt $\log(1/d) \in L^p([0,T])$ fuer alle $p < \infty$
und alle $T > 0$, und die LDI-Bedingung ist erfuellt mit
\[
  \int_0^T \log\frac{1}{d(t)}\,\mathrm{d}t
  \;\le\;
  \lambda T + \log(1/C)\cdot T.
\]
Insbesondere gilt STC mit jedem Exponenten $\alpha > 0$, und
das LDI-Integral ist linear in~$T$ beschraenkt.
\end{lemma}

\begin{proof}
Aus $d(t) \le C\,e^{-\lambda t}$ folgt
$\log(1/d(t)) \ge \lambda t - \log C$, und ebenso
\[
  \log\frac{1}{d(t)}
  \;\le\;
  \log\frac{1}{C\,e^{-\lambda t}}
  \;=\;
  \lambda t + \log\frac{1}{C}.
\]
Daher
\[
  \int_0^T \log\frac{1}{d(t)}\,\mathrm{d}t
  \;\le\;
  \int_0^T \bigl(\lambda t + \log(1/C)\bigr)\,\mathrm{d}t
  \;=\;
  \frac{\lambda T^2}{2} + \log(1/C)\cdot T.
\]
Fuer die LDI-Bedingung \eqref{eq:LDI}, mit $u' \in L^2$:
\[
  \int_0^T \|u'(t)\|\,\log\frac{R}{d(t)}\,\mathrm{d}t
  \;\le\;
  \|u'\|_{L^2}\,
  \Bigl(\int_0^T \bigl(\log(R/C) + \lambda t\bigr)^2\,\mathrm{d}t
  \Bigr)^{1/2}
  \;<\; \infty
\]
nach der H\"olderschen Ungleichung, da der rechte Faktor ein
Polynom in~$T$ ist. Die $L^p$-Integrierbarkeit
$\log(1/d) \in L^p([0,T])$ fuer alle $p < \infty$ folgt, da
$\log(1/d(t)) \le \lambda t + \log(1/C)$, was fuer jedes
endliche~$T$ in $L^\infty([0,T])$ liegt.
\end{proof}

\begin{remark}[Exponentielles Tracking als staerkstes STC-Regime]
\label{rem:exp-tracking-strongest}
Lemma~\ref{lem:stc-tracking} zeigt, dass exponentielles Tracking
das ,,einfachste'' dynamische Regime fuer LDI ist: Die Log-Distanz
ist durch eine lineare Funktion der Zeit beschraenkt, was LDI
automatisch fuer jede Trajektorie mit endlicher Energie
($u' \in L^2$) macht. Dies ist das Regime, das fuer starke
Loesungen der 3D-Navier-Stokes-Gleichungen bedingt auf globale
Regularitaet erwartet wird (Beispiel~\ref{ex:NS}). Der
nicht-triviale Fall tritt ein, wenn exponentielles Tracking
versagt -- z.B.\ nahe potentiellem Blow-up -- wo die
schwaechere STC-Bedingung (polynomieller Sublevel-Abfall) oder
Variante~A (BV-Log-Distanz) die relevanten Werkzeuge werden.
\end{remark}


% ==================================================================
\section{Anwendungen auf dissipative PDEs}
\label{sec:applications}
% ==================================================================

Wir illustrieren, wie der abstrakte Rahmen auf konkrete
PDE-Situationen anwendbar ist.  In jedem Fall besitzt der Semiflow
einen kompakten globalen Attraktor, und die Frage ist, ob TLL und
LDI gelten.

% ------------------------------------------------------------------
\subsection{3D-Navier-Stokes auf $\mathbb{T}^3$}
% ------------------------------------------------------------------

\begin{example}[Navier-Stokes-Gleichungen]
\label{ex:NS}
Sei $H = L^2(\mathbb{T}^3;\,\mathbb{R}^3)$ (divergenzfrei),
$\mathcal{A}$ der globale Attraktor (im Sinne von Temam~\cite{Temam1997}) mit glatter
Antriebskraft~$f$.  Nach Temam~\cite{Temam1997}:
\begin{itemize}
  \item $\mathcal{A}$ ist kompakt in $H^1$ (also auch in $H$).
  \item Jedes $a \in \mathcal{A}$ ist $C^\infty$.
  \item Exponentielles Tracking gilt \emph{bedingt auf globale
    Regularitaet}: Fuer starke Loesungen, die auf $[0,\infty)$
    existieren, gilt $d(t) \le C_0\,e^{-\lambda t}$ fuer $t \ge t_0$.
    (Fuer schwache Leray-Hopf-Loesungen ist nur algebraische
    Attraktion in $L^2$ unbedingt bekannt.)
\end{itemize}

Der Attraktor hat endliche fraktale Dimension
$d_f(\mathcal{A}) \le c_1\,G^{2/3}(1+\log G)^{1/3}$
(wobei $G$ die Grashof-Zahl ist;
vgl.~\cite{ConstantinFoiasTemam1985}), hat aber generisch
Reichweite~$0$ -- was Federers Theorem ausschliesst.

\medskip\noindent
\textbf{Status von TLL:}\; Offen.  Die Squeezing-Eigenschaft des
Navier-Stokes-Semiflows (vgl.~\cite{FoiasManleyRosaTemam2001},
Kapitel~14) liefert exponentielle Kontraktion hochfrequenter
Differenzen auf dem Attraktor.  Kombiniert mit
Man\'e-Projektor-Abschaetzungen deutet dies darauf hin, dass die
Naechstpunkt-Projektion entlang von Trajektorien hoechstens einen
log-Lipschitz-Modul besitzt -- ein rigoroser Beweis erfordert jedoch
quantitative Kontrolle der Projektionsstabilitaet im
Hochmoden-Komplement $Q_N H$.

\medskip\noindent
\textbf{Status von LDI:}\; Folgt aus exponentiellem Tracking
auf $[t_0, T]$ fuer jedes feste $T < \infty$
(Bemerkung~\ref{rem:STC-sources}).  Die kritische Frage betrifft das
Blow-up-Regime $T \nearrow T^*$, wo $\|u'\| \to \infty$.
Die BV-Schranke \eqref{eq:BV-bound} liefert Annahme~G$'$ aus
\cite{FST-NS2026}, die globale Regularitaet ueber das
Attraktorabstands-Argument impliziert.

\medskip\noindent
\textbf{Verbindung zum Log-Dirichlet-Quotienten:}\;
Das Begleitpaper~\cite{FST-NS2026} fuehrt den
\emph{Log-Dirichlet-Quotienten}
$Q_{\mathrm{LD}}(t) = \|\nabla(u-v)\|^2 / \|u-v\|^2$ ein und zeigt,
dass dessen Integrierbarkeit die Gronwall-Abschaetzung von
exponentiellem zu algebraischem Wachstum transformiert. Die
LDI-Bedingung der vorliegenden Arbeit liefert genau den Mechanismus,
der $Q_{\mathrm{LD}}$-Integrierbarkeit sicherstellt: Die BV-Kettenregel
kontrolliert die Komposition $\log d \circ u$, was den
Log-Dirichlet-Quotienten von oben beschraenkt.
\end{example}

% ------------------------------------------------------------------
\subsection{Reaktions-Diffusions-Systeme}
% ------------------------------------------------------------------

\begin{example}[Reaktions-Diffusion auf beschraenkten Gebieten]
\label{ex:RD}
Betrachte $\partial_t u = D\Delta u + f(u)$ auf einem beschraenkten
Gebiet $\Omega \subset \mathbb{R}^n$ mit Dirichlet- oder
Neumann-Randbedingungen, wobei $D > 0$ eine Diffusionsmatrix ist und
$f$ eine Dissipativitaetsbedingung
$\langle f(u), u \rangle \le C - c|u|^2$ erfuellt.  Der Semiflow
besitzt einen kompakten globalen Attraktor
$\mathcal{A} \subset L^2(\Omega)$ mit endlicher fraktaler
Dimension~\cite{Temam1997}.

Fuer skalare Gleichungen ($u \in \mathbb{R}$) in einer Raumdimension
mit Gradienten-Nichtlinearitaet $f(u) = -V'(u)$ besteht der
Attraktor aus Gleichgewichten und heteroklinen Orbits, und die
Naechstpunkt-Projektion ist Lipschitz (tatsaechlich hat der Attraktor
eine Lipschitz-Graph-Struktur als Inertialmannigfaltigkeit;
vgl.~\cite{MalletParetSell1988}).  TLL und LDI gelten trivial.

Fuer Systeme in hoeheren Dimensionen kann der Attraktor
nicht-ganzzahlige fraktale Dimension und Reichweite null haben.
Der log-Lipschitz-Rahmen bietet dann einen Weg zu BV-Selektionen,
den Federers Theorem nicht bereitstellt.  Die Spektrallueckenbedingung
fuer Inertialmannigfaltigkeiten ist fuer viele
Reaktions-Diffusions-Systeme erfuellt~\cite{EdenFoiasNicolaenkoTemam1994},
was die Attraktorgeometrie regulaerer macht als im
Navier-Stokes-Fall.
\end{example}

% ------------------------------------------------------------------
\subsection{Gedaempfte Wellengleichungen}
% ------------------------------------------------------------------

\begin{example}[Gedaempfte semilineare Wellengleichung]
\label{ex:wave}
Betrachte $\partial_{tt} u + \gamma\,\partial_t u = \Delta u + f(u)$
auf $\Omega \subset \mathbb{R}^3$ mit $\gamma > 0$ und subkritischer
Nichtlinearitaet.  Der Semiflow auf $H^1_0 \times L^2$ besitzt einen
kompakten globalen Attraktor mit endlicher fraktaler Dimension
(vgl.~\cite{BabinVishik1992}).  Die exponentielle Tracking-Eigenschaft
gilt mit Rate $\lambda = \gamma/2$ (Spektralluecke des gedaempften
Wellenoperators; vgl.~\cite{BabinVishik1992}, Kapitel~V).

Die Attraktorgeometrie liegt zwischen Reaktions-Diffusion
(wo Inertialmannigfaltigkeiten oft existieren) und Navier-Stokes
(wo sie nicht existieren).  Der log-Lipschitz-Rahmen bietet einen
\emph{einheitlichen} Zugang, der auf alle drei Situationen
anwendbar ist, mit TLL als einzigem geometrischen Input.
\end{example}


% ==================================================================
\section{Diskussion}
\label{sec:discussion}
% ==================================================================

\subsection{Die neue mathematische Schnittstelle}

Der zentrale Beitrag dieser Arbeit ist die Identifikation der
\emph{kritischen Regularitaetsschnittstelle} fuer BV-Selektionen
auf fraktalen Attraktoren.  Die Schnittstelle wird durch eine
einzige Frage charakterisiert:

\begin{center}
\fbox{\parbox{0.85\textwidth}{%
\textbf{Kernfrage.}\;
Erfuellt fuer einen dissipativen Semiflow mit kompaktem Attraktor
$\mathcal{A} \subset H$ die Naechstpunkt-Projektion die
trajektorienweise log-Lipschitz-Kontrolle \eqref{eq:TLL}?
}}
\end{center}

\noindent
Dies ist eine Frage ueber die \emph{Regularitaet des Attraktors
aus Sicht der Trajektorie}, nicht ueber die globale Geometrie
von~$\mathcal{A}$.  Die Unterscheidung ist entscheidend: Globale
geometrische Bedingungen (positive Reichweite, Inertialmannigfaltigkeit)
erzwingen Struktur auf dem \emph{gesamten} Attraktor, waehrend TLL
Regularitaet nur in Richtung der dynamischen Annaeherung erfordert.

\subsection{Bezug zu bestehenden Regularitaetstheorien}

\begin{center}
\begin{tabular}{lllll}
\hline
\textbf{Bedingung} & \textbf{Projektion} & \textbf{Komposition}
  & \textbf{Fraktal?} & \textbf{Quelle} \\
\hline
Positive Reichweite $> 0$
  & Lipschitz-1
  & AC $\Rightarrow$ AC
  & Nein
  & Federer~\cite{Federer1959} \\
Prox-regulaer
  & lokal Lipschitz
  & AC $\Rightarrow$ AC (lokal)
  & Nein
  & Poliquin et al.~\cite{PoliquinRockafellarThibault2000} \\
Hausdorff-BV-Mengenabb.
  & BV-Selektion existiert
  & ---
  & Ja
  & Chistyakov~\cite{Chistyakov2004} \\
Man\'e-Projektor
  & H\"older${}^{-1}$
  & AC $\Rightarrow$ H\"older
  & Ja
  & Man\'e~\cite{Mane1981} \\
\textbf{TLL + LDI}
  & \textbf{log-Lipschitz}
  & \textbf{AC $\Rightarrow$ BV}
  & \textbf{Ja}
  & \textbf{Diese Arbeit} \\
\hline
\end{tabular}
\end{center}

\noindent
Der TLL+LDI-Rahmen ist der einzige Eintrag, der gleichzeitig
(i)~fraktale Geometrie beruecksichtigt (nicht-ganzzahlige Dimension,
Reichweite~$0$), (ii)~BV-Regularitaet liefert (hinreichend fuer
Gronwall-artige Argumente) und (iii)~als Trajektorien-Bedingung
formuliert ist (unabhaengig von der globalen Attraktorstruktur).

\subsection{Einschraenkungen}

\begin{itemize}
  \item \textbf{Selektionsabhaengigkeit.}\;
    Die TLL-Bedingung ist fuer eine \emph{gegebene} messbare
    Naechstpunkt-Selektion~$P$ formuliert.  Auf nichtkonvexen
    kompakten Mengen ist die Naechstpunkt-Abbildung generisch
    mehrwertig, und verschiedene Selektionen koennen unterschiedliche
    Regularitaet haben.  Der Rahmen erfordert die Existenz
    \emph{mindestens einer} Selektion, die TLL erfuellt; er
    behauptet nicht, dass jede messbare Selektion dies tut.
    In der Praxis ist die ,,beste'' Selektion (z.B.\ diejenige,
    die die metrische Ableitung minimiert) der natuerliche Kandidat.
  \item \textbf{TLL ist offen fuer 3D-NS.}\;
    Die Squeezing-Eigenschaft und Man\'e-Projektor-Abschaetzungen
    deuten auf log-Lipschitz-Stabilitaet der Projektion hin, aber ein
    rigoroser Beweis erfordert quantitative Schranken fuer das
    Hochmoden-Projektionskomplement, die derzeit nicht verfuegbar sind.
  \item \textbf{LDI ist nur unter Blow-up nicht-trivial.}\;
    Fuer glatte Loesungen auf kompakten Zeitintervallen gilt LDI
    automatisch.  Die Bedingung wird erst im Grenzfall
    $T \nearrow T^*$ (Blow-up-Zeit) kritisch, wo $\|u'\| \to \infty$.
  \item \textbf{BV, nicht AC.}\;
    Der Rahmen liefert BV-Regularitaet, die fuer
    Stieltjes-Integral-basierte Gronwall-Argumente genuegt, aber nicht
    fuer punktweise Differentialungleichungen.  Ein Upgrade von BV auf
    AC wuerde Lipschitz-Projektion (positive Reichweite) erfordern --
    genau die Bedingung, die wir zu vermeiden suchen.
\end{itemize}


% ==================================================================
\section{Schlussfolgerung}
% ==================================================================

Wir haben einen Regularitaetsrahmen fuer Naechstpunkt-Projektionen
auf kompakte Attraktoren eingefuehrt, der die kritische Luecke
zwischen H\"older (unzureichend fuer BV) und Lipschitz (unvereinbar
mit fraktaler Geometrie) besetzt.  Die trajektorienweise
log-Lipschitz-Bedingung (TLL) kontrolliert den Projektionsmodul auf
der Attraktorabstandsskala, und die
Log-Distanz-Integrierbarkeitsbedingung (LDI) stellt sicher, dass die
resultierenden BV-Summen konvergieren.

Der Rahmen reduziert bedingte Regularitaetskriterien fuer dissipative
PDEs -- einschliesslich der 3D-Navier-Stokes-Gleichungen -- auf eine
einzige geometrisch-dynamische Frage: Hat die Naechstpunkt-Projektion
hoechstens logarithmisches Blow-up entlang dissipativer Trajektorien?
Diese Frage ist unabhaengig von der spezifischen PDE und haengt nur
vom Zusammenspiel zwischen Attraktorgeometrie und den
Kontraktionseigenschaften des Semiflows ab.


% ==================================================================
\section{Zeitgemittelte Selektion und BV-Regularitaet}
\label{sec:time-averaged}
% ==================================================================

Die zentrale Obstruktion zur Schliessung der Regularitaetsluecke ist
die BV-Kontrolle der Naechstpunkt-Selektion
$v(t) \in \arg\min_{a \in \mathcal{A}} \|u(t) - a\|$.  Wir schlagen
vor, die punktweise Selektion durch eine \emph{zeitgemittelte}
Selektion zu ersetzen, die automatisch BV-regulaer ist.

\begin{definition}[Moreau--Yosida-regularisierte Selektion]
\label{def:my-selection}
Fuer $\varepsilon > 0$ definiere die \emph{zeitgemittelte
Naechstpunkt-Selektion}
\[
  v_\varepsilon(t)
  \;:=\;
  \arg\min_{a \in \mathcal{A}}
  \int_{t-\varepsilon}^{t+\varepsilon}
  \|u(s) - a\|^2 \, ds.
\]
\end{definition}

\begin{proposition}[Automatische BV-Regularitaet]
\label{prop:bv-regularity}
Unter den Standardannahmen (endlich-dimensionaler Attraktor,
Energiedissipationsschranke) erfuellt die regularisierte
Selektion~$v_\varepsilon$:
\begin{enumerate}[label=(\roman*)]
\item \emph{Eindeutigkeit:}  $v_\varepsilon(t)$ ist f.ue.\ eindeutig
  (strikte Konvexitaet des zeitintegrierten Kostenfunktionals).
\item \emph{BV-Schranke:}
  $\Var(v_\varepsilon;\,[0,T])
   \;\le\;
   C_\varepsilon
   \int_0^T \|\partial_t u(s)\|\,ds
   \;<\; \infty$.
\item \emph{TLL-Vererbung:}  Falls~$u$ TLL mit Konstante
  $C_{\mathrm{TLL}}$ erfuellt, so erfuellt $v_\varepsilon$ TLL mit
  Konstante $C_{\mathrm{TLL}} + O(\varepsilon)$.
\item \emph{Gronwall-Kompatibilitaet:}  Die Gronwall-Abschaetzungen
  des Hauptresultats (Satz~\ref{thm:BV-chain}) bleiben fuer
  $v_\varepsilon$ gueltig mit einem zusaetzlichen Fehlerterm
  $O(\varepsilon\,\|\partial_t u\|_{L^2})$.
\end{enumerate}
\end{proposition}

\begin{proof}[Beweisskizze]
\emph{(i)}\; Die Abbildung
$a \mapsto \int_{t-\varepsilon}^{t+\varepsilon} \|u(s)-a\|^2\,ds$
ist strikt konvex: Der Integrand ist strikt konvex fuer f.ue.~$s$
(da $u(s) \neq a$ fuer f.ue.~$s$ und $a \in \mathcal{A}$), und
Integration erhaelt strikte Konvexitaet.

\emph{(ii)}\; Die Erste-Variations-Bedingung fuer den Minimierer
liefert
\[
  \|v_\varepsilon(t+h) - v_\varepsilon(t)\|
  \;\le\;
  \frac{1}{2\varepsilon}
  \int_{t-\varepsilon}^{t+\varepsilon}
  \|u(s+h) - u(s)\|\,ds
  \;\le\;
  \frac{h}{2\varepsilon}
  \int_{t-\varepsilon}^{t+\varepsilon}
  \|\partial_t u(s)\|\,ds.
\]
Integration ueber $t \in [0,T]$ und Fubini ergibt die BV-Schranke
mit $C_\varepsilon = T/(2\varepsilon)$.

\emph{(iii)}\; Die TLL-Konstante verschlechtert sich um das
Zeitmittelungsfenster:
\[
  |\log d_\varepsilon(t+h) - \log d_\varepsilon(t)|
  \;\le\;
  C_{\mathrm{TLL}}\,|\log|h||
  \;+\;
  C\,\varepsilon/|h|,
\]
wobei der zweite Term aus der Diskrepanz zwischen punktweisem und
gemitteltem Abstand entsteht.

\emph{(iv)}\; Die Gronwall-Abschaetzung nimmt einen zusaetzlichen
Forcingterm aus der Zeitmittelung auf:
\[
  \frac{d}{dt} H(u \mid v_\varepsilon)
  \;\le\;
  -\alpha\,H + \beta + \gamma\,\varepsilon,
\]
wobei $\gamma$ von $\|\partial_t u\|$ abhaengt, aber nicht von~$H$.
Das Standard-Gronwall-Lemma ist anwendbar; der zusaetzliche
$\gamma\varepsilon$-Term verschwindet im Limes
$\varepsilon \to 0$.
\end{proof}

\begin{remark}[Grenzuebergang $\varepsilon \to 0$]
\label{rem:epsilon-limit}
Die zentrale Frage ist, ob die BV-Schranke den Grenzuebergang
$\varepsilon \to 0$ ueberlebt:
\begin{itemize}
\item Falls\; $\sup_\varepsilon \Var(v_\varepsilon;\,[0,T]) < \infty$,
  konvergiert eine Teilfolge $v_{\varepsilon_k}$ in~$L^1$ gegen eine
  BV-Selektion~$v_0$, die die \emph{kanonische} Naechstpunkt-Selektion
  ist.
\item Die gleichmaessige BV-Schranke gilt genau dann, wenn der
  Attraktor $\mathcal{A}$ eine ,,zeitgemittelte Reichweite''
  $\mathrm{reach}_{\mathrm{avg}}(\mathcal{A}) > 0$ besitzt, die
  echt schwaecher ist als positive Reichweite: Sie erlaubt isolierte
  Spitzen vom Mass null.
\end{itemize}
Dies ist die praezise Formulierung der ,,BV-Regularitaetsluecke'' und
verbindet sich mit:
\begin{itemize}
\item \emph{Moreau--Yosida-Regularisierung} in der konvexen Analysis
  (die zeitgemittelten Kosten sind die Moreau-Enveloppe),
\item \emph{Viskoser Selektion} in Gradientensystemen (die Zeitmittelung
  wirkt als kuenstliche Viskositaet),
\item \emph{Thermodynamischer Selektion} im Dark-Energy-Begleitpaper
  (die IR-Mittelung von~$\Phi[\rho]$).
\end{itemize}
\end{remark}


% ==================================================================
\section{Offene Richtungen}
\label{sec:open-directions}
% ==================================================================

Der TLL+LDI-Rahmen eroeffnet mehrere konkrete Forschungsrichtungen,
die thematisch gruppiert sind.

% ------------------------------------------------------------------
\subsection{Pilotsysteme zur Verifikation von TLL und LDI}
% ------------------------------------------------------------------

\begin{remark}[Pilotsysteme zur Verifikation von TLL und LDI]
\label{rem:pilot-verification}
Drei Kandidatensysteme zur Pilotverifikation der TLL- und LDI-Bedingungen:
\begin{enumerate}
\item \emph{Reaktions-Diffusion (2D):} Glatte Nichtlinearitaet, gut verstandene Attraktoren. TLL und LDI gelten erwartungsgemaess aufgrund der parabolischen Glaettung.
\item \emph{Hyperviskose NSE ($(-\Delta)^\beta$, $\beta \geq 5/4$):} Inertialmannigfaltigkeiten existieren (Mallet-Paret--Sell, 1988), was TLL garantiert. Dient als Brueckensystem zwischen dem handhabbaren Reaktions-Diffusions-Fall und den vollen NS ($\beta = 1$).
\item \emph{Kuramoto--Sivashinsky (1D):} Deterministisches Chaos ohne Singularitaeten. Der Attraktor ist endlich-dimensional und glatt, was TLL+LDI durch direkte Berechnung verifizierbar macht.
\end{enumerate}
\end{remark}

\begin{remark}[Ginzburg--Landau als TLL/LDI-Bruecke]
\label{rem:ginzburg-landau}
Die komplexe Ginzburg--Landau-Gleichung $\partial_t u = (1+i\alpha)\Delta u + u - (1+i\beta)|u|^2 u$ besitzt fuer bestimmte Parameterbereiche eine endlich-dimensionale Inertialmannigfaltigkeit. Vollstaendige Verifikation von TLL+LDI auf diesem System, gefolgt von einem Lifting-Argument auf 2D-NS (via Parameterfortsetzung in der Nichtlinearitaet), wuerde einen konkreten Weg zum vollen NS-Regularitaetsproblem liefern.
\end{remark}

% ------------------------------------------------------------------
\subsection{Mathematische Werkzeuge fuer TLL und LDI}
% ------------------------------------------------------------------

\begin{remark}[Brueckenlemma: Squeezing zu TLL]
\label{rem:bridge-squeezing}
Die Squeezing-Eigenschaft (Foias--Temam) besagt, dass hohe Moden schneller kontrahieren als niedrige Moden divergieren. Kombiniert mit dem Theorem der bestimmenden Moden impliziert dies: $d(u(t), \mathcal{A}) \leq C e^{-\gamma_N t} d(u(0), \mathcal{A}) + C' \sup_{s \leq t} |P_N u(s) - P_N v(s)|$ fuer ein $v \in \mathcal{A}$. Logarithmierung ergibt $\log(1/d(t)) \geq \gamma_N t - C''$, was die dem TLL zugrundeliegende Log-Modul-Schranke ist.
\end{remark}

\begin{remark}[Verdopplungs- und Assouad-Struktur]
\label{rem:doubling-assouad}
TLL kann in Begriffen der Assouad-Dimension der Attraktortrajektorie umformuliert werden: Wenn die Trajektorie $\{u(t) : t \in [0,T]\}$ Assouad-Dimension $< 1$ hat, dann hat der ,,approximative Tangentialkegel'' einen Oeffnungswinkel beschraenkt durch $C/|\log r|$, was aequivalent zu TLL mit Rate $1/|\log r|$ ist. Dies verbindet die Regularitaetstheorie mit der geometrischen Masstheorie. Neuere Strukturresultate ueber Mengen positiver Reichweite~\cite{Lytchak2024} liefern zusaetzliche Werkzeuge zum Verstaendnis der Geometrie von Attraktorumgebungen.
\end{remark}

\begin{remark}[LDI als Schwanzintegral-Aequivalenz]
\label{rem:ldi-tail}
Die LDI-Bedingung $\int_0^T \log(1/d(t))\,\mathrm{d}t < \infty$ ist aequivalent zur Schwanzintegral-Bedingung $\int_0^R \frac{1}{s} \left( \int_{\{0 < d \leq s\}} \|u'(t)\|\,\mathrm{d}t \right) \mathrm{d}s < \infty$, die die Integrierbarkeit von $\log(1/d)$ mit der ,,Geschwindigkeit'' der Trajektorie nahe dem Attraktor in Beziehung setzt. Diese Formulierung ist fuer Energieabschaetzungen besser geeignet.
\end{remark}

\begin{remark}[Klaerung zu Proposition~\ref{prop:variant-B}]
\label{rem:prop-variant-B-clarification}
Die BV $\to$ $L^\infty$-Einbettung in Proposition~\ref{prop:variant-B} erfordert die zusaetzliche Bedingung $d(t) > 0$ f.\"u.: Die Trajektorie darf den Attraktor nicht auf einer Menge positiven Masses beruehren. Dies ist fuer Leray--Hopf-Loesungen automatisch erfuellt (die analytisch im Raum sind und daher nicht mit einem Attraktorelement auf einem Intervall uebereinstimmen koennen, es sei denn sie sind dieses Element).
\end{remark}

% ------------------------------------------------------------------
\subsection{Stochastische Erweiterungen und fortgeschrittene Werkzeuge}
% ------------------------------------------------------------------

\begin{remark}[Stochastische Stabilisierung]
\label{rem:stochastic-stabilisation}
Hinzufuegen von Feller-artigem Rauschen $\sigma\,\mathrm{d}W$ zu den Navier--Stokes-Gleichungen kann TLL ,,fast sicher'' erzwingen: Das Rauschen verhindert, dass die Trajektorie sich dem Attraktor tangential naehert (was log-Lipschitz verletzen wuerde), und die Feller-Eigenschaft stellt sicher, dass die Uebergangshalbgruppe beschraenkte Funktionen auf stetige abbildet, was die fuer das Selektionsprinzip benoetigte Regularitaet liefert.
\end{remark}

\begin{remark}[Hochmoden-Stabilitaetsvermutung]
\label{rem:high-mode-stability}
Eine hinreichende Bedingung fuer TLL im NS-Kontext: Wenn die hohen Fourier-Moden $Q_N u = (I - P_N) u$ die Bedingung $\|Q_N u(t)\| \leq C N^{-\alpha}$ gleichmaessig auf dem Attraktor fuer ein $\alpha > 1/2$ erfuellen, dann gilt $d(t) \leq C' N^{-\alpha}$ und TLL folgt aus der Abschaetzung der bestimmenden Moden. Diese ,,Hochmoden-Stabilitaetsvermutung'' ist schwaecher als die Existenz einer Inertialmannigfaltigkeit (die $\alpha > 1$ erfordert) und koennte durch Galerkin-Energieabschaetzungen verifizierbar sein. Neuere scharfe Dimensionsabschaetzungen fuer NS-artige Attraktoren~\cite{IlyinKalantarovZelik2025} liefern quantitativen Input fuer solche Schranken.
\end{remark}

\begin{remark}[Perspektive der Variationsanalysis]
\label{rem:variational-analysis}
Die Theorie der metrischen Regularitaet und Variationsanalysis (Rockafellar--Wets~\cite{RockafellarWets1998}, Dontchev--Rockafellar~\cite{DontchevRockafellar2009}) bietet eine natuerliche Heimat fuer die TLL- und LDI-Bedingungen: TLL steht in Beziehung zur Aubin-Eigenschaft der Attraktorprojektion, und LDI ist eine Calmness-Bedingung fuer die Distanzfunktion. Das vollstaendige Werkzeugset der Variationsanalysis -- einschliesslich Ekelands Variationsprinzip, des Mordukhovich-Kriteriums und der Attouch--Th\'era-Dualitaet -- ist anwendbar und in diesem Kontext wenig erforscht.
\end{remark}

% ------------------------------------------------------------------
\subsection{Weitere NS-spezifische Forschungsrichtungen}
% ------------------------------------------------------------------

\begin{remark}[Weitere Forschungsrichtungen]
\label{rem:further-ns-logdist}
\begin{enumerate}
\item \emph{Selektionsregularitaet ueber Trajektorienattraktoren:} Formulierung der minimalen metrischen Ableitungsselektion im Chepyzhov--Vishik-Trajektorienattraktor-Rahmen, wo Nichteindeutigkeit durch Betrachtung ganzer Trajektorien statt Anfangsdaten behandelt wird.
\item \emph{Quantitative Tracking-Lemmata:} Herleitung von Massabschaetzungen $|\{t : d(t) \leq \varepsilon\}| \leq C \varepsilon^\alpha$ ueber Energie-/Absorbierende-Kugel-Argumente und ,,Kein-klebriger-Kontakt''-Bedingungen.
\item \emph{Ekelands Variationsprinzip:} Nutzung von Ekelands Prinzip zur Konstruktion von ,,Fast-Minimierern'' mit kontrollierter Variation, wodurch die Notwendigkeit positiver Reichweite von $\mathcal{A}$ umgangen wird.
\item \emph{Mengenwertige BV-Selektion:} Bestimmung, welche dynamischen Regularitaetsbedingungen an den Fluss beschraenkte Hausdorff-Variation der Minimierermenge $\{u \in \mathcal{A} : d(u, u(t)) = d_{\mathcal{A}}(u(t))\}$ erzwingen.
\item \emph{Approximative Prox-Regularitaet:} Herstellung einer modifizierten Reichweite nach Entfernung einer meageren Menge (Lebesgue-Null): Der Attraktor kann auf einer vernachlaessigbaren Menge keine positive Reichweite haben, sie aber generisch beibehalten. Resultate von Lytchak~\cite{Lytchak2024} ueber die Struktur von Mengen positiver Reichweite legen nahe, dass eine solche Zerlegung durchfuehrbar sein koennte.
\item \emph{Schwache Selektionsfunktionen:} Ersetzung starker Lipschitz-Projektionen durch Selektionsfunktionen beschraenkter Hausdorff-Variation, die unter schwaecheren Regularitaetsannahmen existieren.
\item \emph{Geometrisches Euler-System:} Konstruktion einer kohaerenten Hierarchie endlicher Galerkin-Approximationen $\mathcal{A}_N \to \mathcal{A}$ als ,,Reichweiten-Surrogat'': Selbst wenn $\mathrm{reach}(\mathcal{A}) = 0$, haben die Galerkin-Approximationen $\mathcal{A}_N$ positive Reichweite und konvergieren im Hausdorff-Abstand.
\end{enumerate}
\end{remark}

% ------------------------------------------------------------------
\subsection{Paper-uebergreifende Verbindungen}
% ------------------------------------------------------------------

\begin{remark}[DFC aus Turbulenz als topologischer Trichter fuer TLL]
\label{rem:dfc-funnel}
Die Detaillierte Flussbedingung (DFC) aus dem Begleitpaper zur Turbulenz~\cite{FST-NS2026} impliziert, dass Energie ,,abwaerts'' durch den Inertialbereich fliesst. Uebersetzt in die Attraktorgeometrie bedeutet dies, dass Trajektorien sich $\mathcal{A}$ durch einen topologischen Trichter naehern: Die DFC beschraenkt die Annaeherungsrichtung und erzwingt eine log-Lipschitz-Rate. Dies liefert einen physikalischen Mechanismus fuer TLL, der unabhaengig von den abstrakten funktionalanalytischen Argumenten ist.
\end{remark}


% ===================================================================
% KI-Offenlegung
% ===================================================================
\subsection*{KI-Offenlegung}
Bei der Erstellung dieses Manuskripts wurden KI-gestuetzte Werkzeuge
(Claude, Anthropic) in unterstuetzender Funktion eingesetzt,
insbesondere fuer Literaturrecherche, Textstrukturierung und
Formulierungshilfe.  Die konzeptionelle Gestaltung, die
wissenschaftliche Argumentation und die Verantwortung fuer den
gesamten Inhalt liegen ausschliesslich beim Autor.


\begin{thebibliography}{99}

\bibitem{AmbrosioGigliSavare2008}
L.~Ambrosio, N.~Gigli, and G.~Savar\'e,
\emph{Gradient Flows in Metric Spaces and in the Space of
Probability Measures},
2nd ed., Lectures in Mathematics ETH Z\"urich,
Birkh\"auser, Basel, 2008.

\bibitem{AmbrosioFuscoPallara2000}
L.~Ambrosio, N.~Fusco, and D.~Pallara,
\emph{Functions of Bounded Variation and Free Discontinuity Problems},
Oxford Mathematical Monographs, Clarendon Press, Oxford, 2000.

\bibitem{BabinVishik1992}
A.~V.~Babin and M.~I.~Vishik,
\emph{Attractors of Evolution Equations},
Studies in Mathematics and its Applications, vol.~25,
North-Holland, Amsterdam, 1992.

\bibitem{BarbuCannone2016}
V.~Barbu and M.~Cannone,
\emph{Locally Lipschitz composition operators in spaces of
functions of bounded variation},
arXiv preprint arXiv:1602.04490, 2016.

\bibitem{Chistyakov2004}
V.~V.~Chistyakov,
\emph{Selections of bounded variation},
J.~Appl.\ Anal.\ \textbf{10} (2004), no.~1, 1--82.

\bibitem{DontchevRockafellar2009}
A.~L.~Dontchev and R.~T.~Rockafellar,
\emph{Implicit Functions and Solution Mappings: A View from
Variational Analysis},
Springer Monographs in Mathematics, Springer, New York, 2009.

\bibitem{ConstantinFoiasTemam1985}
P.~Constantin, C.~Foias, and R.~Temam,
\emph{Attractors representing turbulent flows},
Mem.\ Amer.\ Math.\ Soc.\ \textbf{53} (1985), no.~314.

\bibitem{EdenFoiasNicolaenkoTemam1994}
A.~Eden, C.~Foias, B.~Nicolaenko, and R.~Temam,
\emph{Exponential Attractors for Dissipative Evolution Equations},
Research in Applied Mathematics, vol.~37,
Wiley--Masson, Paris, 1994.

\bibitem{Federer1959}
H.~Federer,
\emph{Curvature measures},
Trans.\ Amer.\ Math.\ Soc.\ \textbf{93} (1959), 418--491.

\bibitem{FoiasManleyRosaTemam2001}
C.~Foias, O.~Manley, R.~Rosa, and R.~Temam,
\emph{Navier--Stokes Equations and Turbulence},
Encyclopedia of Mathematics and its Applications, vol.~83,
Cambridge University Press, 2001.

\bibitem{FST-NS2026}
L.~Geiger,
\emph{A Thermodynamic Obstruction to Finite-Time Blow-Up in the 3D
Navier--Stokes Equations (Conditional Framework)},
Zenodo, 2026.  DOI:\@ to be assigned.

\bibitem{IlyinKalantarovZelik2025}
A.~Ilyin, V.~Kalantarov, and S.~Zelik,
\emph{Attractors and their dimensions for the 3D Fractional
Navier--Stokes--Voigt Equations},
arXiv preprint arXiv:2511.04783, 2025.

\bibitem{Lytchak2024}
A.~Lytchak,
\emph{A note on subsets of positive reach},
Math.\ Nachr.\ \textbf{297} (2024), no.~3, 932--942.

\bibitem{MalletParetSell1988}
J.~Mallet-Paret and G.~R. Sell,
\emph{Inertial manifolds for reaction diffusion equations in higher
space dimensions},
J.~Amer.\ Math.\ Soc.\ \textbf{1} (1988), no.~4, 805--866.

\bibitem{Mane1981}
R.~Man\'e,
\emph{On the dimension of the compact invariant sets of certain
non-linear maps},
Lecture Notes in Mathematics, vol.~898, Springer, Berlin, 1981,
pp.~230--242.

\bibitem{RockafellarWets1998}
R.~T.~Rockafellar and R.~J-B Wets,
\emph{Variational Analysis},
Grundlehren der mathematischen Wissenschaften, vol.~317,
Springer, Berlin, 1998.

\bibitem{PoliquinRockafellarThibault2000}
R.~A.~Poliquin, R.~T.~Rockafellar, and L.~Thibault,
\emph{Local differentiability of distance functions},
Trans.\ Amer.\ Math.\ Soc.\ \textbf{352} (2000), no.~11, 5231--5249.

\bibitem{Saks1937}
S.~Saks,
\emph{Theory of the Integral},
2nd revised ed., Monografie Matematyczne, vol.~7,
G.~E.~Stechert \& Co., New York, 1937.

\bibitem{Temam1997}
R.~Temam,
\emph{Infinite-Dimensional Dynamical Systems in Mechanics and Physics},
2nd ed., Applied Mathematical Sciences, vol.~68,
Springer, New York, 1997.

\end{thebibliography}

\end{document}
